\NormalFont\ShortTitle{Romani}

{\MTB Scrisoarea lui Pavel către

\par }{\MT Romani


\par }\ChapOne{1}{\PP \VerseOne{1}Pavel, slujitor al lui Isus Hristos,
\FTNT{*}{{\FR 1:1 }{\FT{“Hristos” înseamnă “Unsul”.}}}chemat ca apostol, pus deoparte pentru vestea cea bună a lui Dumnezeu,

\VS{2}pe care a făgăduit-o mai înainte prin proorocii Săi în Sfintele Scripturi,

\VS{3}despre Fiul Său, născut din neamul\FTNT{†}{{\FR 1:3 }{\FT{sau, semințe}}} lui David, după trup,

\VS{4}care a fost declarat Fiul lui Dumnezeu cu putere, după Duhul sfințeniei, prin învierea din morți, Isus Hristos, Domnul nostru,

\VS{5}prin care am primit harul și apostolatul pentru ascultarea credinței printre toate neamurile, pentru Numele Lui;

\VS{6}printre care și voi sunteți chemați să aparțineți lui Isus Cristos;

\VS{7}tuturor celor care sunt în Roma, iubiți de Dumnezeu, chemați să fie sfinți: Harul și pacea să vă fie vouă, de la Dumnezeu Tatăl nostru și de la Domnul Isus Hristos.


\par }{\PP \VS{8}Mai întâi, mulțumesc Dumnezeului meu, prin Isus Hristos, pentru voi toți, pentru că credința voastră este vestită în toată lumea.

\VS{9}Căci martor îmi este Dumnezeu, căruia îi slujesc în duhul meu în Buna Vestire a Fiului său, cât de neîncetat fac mereu pomenirea voastră în rugăciunile mele,

\VS{10}cerând, dacă acum, în sfârșit, prin voia lui Dumnezeu, să fiu învrednicit să ajung la voi.

\VS{11}Căci doresc mult să vă văd, ca să vă pot împărtăși vreun dar spiritual, pentru ca voi să fiți întăriți;

\VS{12}adică, ca eu împreună cu voi să fiu încurajat în voi, fiecare dintre noi prin credința celuilalt, atât a voastră, cât și a mea.


\par }{\PP \VS{13}Și nu vreau să nu știți, fraților, că de multe ori am plănuit să vin la voi, dar am fost împiedicat până acum, ca să aduc rod și printre voi, ca și printre celelalte neamuri.

\VS{14}Sunt dator atât grecilor, cât și străinilor, atât celor înțelepți, cât și celor nebuni.

\VS{15}Așadar, atât cât este în mine, sunt nerăbdător să propovăduiesc vestea cea bună și vouă, care sunteți în Roma.


\par }{\PP \VS{16}Căci nu mi-e rușine de vestea cea bună a lui Hristos, pentru că ea este puterea lui Dumnezeu pentru mântuirea fiecăruia care crede, mai întâi a iudeilor și apoi a grecilor.

\VS{17}Căci în ea se descoperă dreptatea lui Dumnezeu, din credință în credință. După cum este scris: “Dar cel neprihănit va trăi prin credință”.
\FTNT{‡}{{\FR 1:17 }{\FT{Habacuc 2:4}}}


\par }{\PP \VS{18}Căci mânia lui Dumnezeu se arată din ceruri împotriva oricărei nelegiuiri și nedreptăți a oamenilor care înăbușă adevărul în nedreptate,

\VS{19}pentru că ceea ce este cunoscut de Dumnezeu se arată în ei, căci Dumnezeu le-a descoperit.

\VS{20}Căci lucrurile invizibile ale Lui, de la crearea lumii, se văd în mod clar, fiind percepute prin lucrurile care sunt făcute, adică puterea Lui veșnică și divinitatea Lui, pentru ca ei să fie fără scuză.

\VS{21}Pentru că, cunoscând pe Dumnezeu, nu L-au slăvit ca Dumnezeu și nu I-au adus mulțumiri, ci au devenit deșerți în raționamentul lor și inima lor fără minte s-a întunecat.


\par }{\PP \VS{22}Și, făcându-se înțelepți, s-au făcut nebuni,

\VS{23}și au schimbat slava Dumnezeului cel neînsuflețit cu asemănarea unui chip de om coruptibil, de păsări, de patrupede și de târâtoare.

\VS{24}De aceea, Dumnezeu i-a și predat, în poftele inimilor lor, la necurăție, pentru ca trupurile lor să fie dezonorate între ei;

\VS{25}care au schimbat adevărul lui Dumnezeu cu minciuna și s-au închinat și au slujit mai degrabă creaturii decât Creatorului, care este binecuvântat în veci. Amin.


\par }{\PP \VS{26}De aceea Dumnezeu i-a lăsat să se lase pradă patimilor josnice. Căci femeile lor au schimbat funcția naturală în ceea ce este împotriva naturii.

\VS{27}La fel și bărbații, părăsind funcția naturală a femeii, s-au aprins în poftele lor unii față de alții, bărbații făcând ceea ce este nepotrivit cu bărbații și primind în ei înșiși pedeapsa cuvenită pentru greșeala lor.

\VS{28}Chiar dacă au refuzat să-L aibă pe Dumnezeu în cunoștința lor, Dumnezeu i-a predat unei minți reprobabile, ca să facă ceea ce nu se cuvine;

\VS{29}fiind plini de toată nedreptatea, de imoralitate sexuală, de răutate, de lăcomie, de răutate; plini de invidie, de crimă, de ceartă, de înșelăciune, de obiceiuri rele, calomniatori în secret,

\VS{30}hulitori, urâcioși față de Dumnezeu, obraznici, aroganți, lăudăroși, inventatori de lucruri rele, neascultători față de părinți,

\VS{31}lipsiți de înțelegere, călcători de legământ, fără afecțiune firească, neiertători, nemiloși;

\VS{32}care, cunoscând rânduiala lui Dumnezeu, că cei care practică astfel de lucruri sunt vrednici de moarte, nu numai că fac același lucru, dar și aprobă pe cei care le practică.



\par }\Chap{2}{\PP \VerseOne{1}De aceea, omule, oricine ai fi tu, care judeci, ești fără prihană. Căci în ceea ce judeci pe altul, te osândești pe tine însuți. Căci tu, care judeci, practici aceleași lucruri.

\VS{2}Noi știm că judecata lui Dumnezeu este conformă cu adevărul împotriva celor care practică astfel de lucruri.

\VS{3}Crezi tu, omule, care judeci pe cei care practică astfel de lucruri și faci la fel, că vei scăpa de judecata lui Dumnezeu?

\VS{4}Sau disprețuiești tu bogăția bunătății, a îngăduinței și a răbdării Lui, fără să știi că bunătatea lui Dumnezeu te conduce la pocăință?

\VS{5}Ci, după împietrirea voastră și inima voastră neîmpăcată, vă strângeți pentru voi înșivă mânie în ziua mâniei, a revelației și a dreptei judecăți a lui Dumnezeu,

\VS{6}care “va răsplăti fiecăruia după faptele sale”.
\FTNT{*}{{\FR 2:6 }{\FT{{\CHAR{Psalmul 62:12}}; Proverbe 24:12}}}

\VS{7}celor care, prin perseverența în binele făcut, caută gloria, onoarea și incoruptibilitatea, viața veșnică;

\VS{8}dar celor care sunt egoiști și nu ascultă de adevăr, ci ascultă de nedreptate, le va fi mânia, indignarea,

\VS{9}asuprirea și chinul asupra oricărui suflet de om care face răul, mai întâi iudeilor și apoi grecilor.


\par }{\PP \VS{10}Dar slava, cinstea și pacea sunt pentru orice om care face binele, mai întâi pentru iudeu și apoi pentru grec.

\VS{11}Căci la Dumnezeu nu este părtinire.

\VS{12}Căci toți cei ce au păcătuit fără lege vor pieri și fără lege. Toți cei care au păcătuit sub Lege vor fi judecați prin Lege.

\VS{13}Căci nu ascultătorii legii sunt cei care sunt drepți înaintea lui Dumnezeu, ci cei care împlinesc legea vor fi îndreptățiți

\VS{14}(pentru că atunci când neamurile care nu au legea fac prin natură lucrurile legii, aceștia, neavând legea, sunt o lege pentru ei înșiși,

\VS{15}prin faptul că arată lucrarea legii scrisă în inimile lor, conștiința lor mărturisind cu ei, iar gândurile lor între ei acuzându-i sau dimpotrivă scuzându-i)

\VS{16}în ziua în care Dumnezeu va judeca tainele oamenilor, potrivit cu Buna Vestire a mea, prin Isus Hristos.


\par }{\PP \VS{17}Tu porți numele de iudeu, te sprijini pe Lege, te lauzi în Dumnezeu,

\VS{18}cunoști voia Lui și aprobi lucrurile bune, fiind învățat din Lege,

\VS{19}și ești încredințat că tu însuți ești călăuză orbilor, lumină pentru cei ce sunt în întuneric,

\VS{20}îndrumător al celor nebuni, învățător al copiilor, având în Lege forma cunoștinței și a adevărului.

\VS{21}Așadar, tu, care înveți pe altul, nu te înveți și pe tine însuți? Tu, care propovăduiești că omul nu trebuie să fure, tu furi?

\VS{22}Tu, care spui că un om nu trebuie să comită adulter, comiți tu adulter? Voi, care aveți în ură idolii, jefuiți templele?

\VS{23}Voi, care vă lăudați cu legea, îl dezonorați pe Dumnezeu prin nerespectarea legii?

\VS{24}Căci “numele lui Dumnezeu este hulit printre neamuri din cauza voastră”\FTNT{†}{{\FR 2:24 }{\FT{{\CHAR{Isaia 52:5}}; Ezechiel 36:22}}}, așa cum este scris.

\VS{25}Căci, într-adevăr, circumcizia este de folos, dacă ești împlinitor al legii, dar dacă ești călcător al legii, circumcizia ta a devenit necircumcizie.

\VS{26}Așadar, dacă cel netăiat împrejur ține rânduielile Legii, nu cumva netăierea lui împrejur nu va fi socotită ca fiind circumcizie?

\VS{27}Nu te vor judeca oare pe tine, care cu litera și cu circumcizia ești un călcător de lege, cei care sunt fizic netăiați împrejur, dar care împlinesc legea?

\VS{28}Căci nu este iudeu cel care este unul în afară, nici circumcizia care este exterioară în carne;

\VS{29}ci este iudeu cel care este unul înăuntru, iar circumcizia este cea a inimii, în duh, nu în literă; a cărui laudă nu vine de la oameni, ci de la Dumnezeu.



\par }\Chap{3}{\PP \VerseOne{1}Atunci ce avantaj are evreul? Sau care este folosul circumciziei?

\VS{2}Mult în toate privințele! Pentru că, mai întâi de toate, li s-au încredințat revelațiile lui Dumnezeu.

\VS{3}Căci ce s-ar întâmpla dacă unii ar fi fără credință? Oare lipsa lor de credință va anula credincioșia lui Dumnezeu?

\VS{4}Să nu se întâmple niciodată! Da, să fie găsit Dumnezeu adevărat, dar orice om mincinos. După cum este scris,

\par }{\Q “pentru ca să fiți îndreptățiți în cuvintele voastre,

\par }{\QB și ar putea să prevaleze atunci când veți veni la judecată.”
\FTNT{*}{{\FR 3:4 }{\FT{{\CHAR{Psalmul 51:4}}}}}


\par }{\PP \VS{5}Dar dacă neprihănirea noastră face să se recunoască neprihănirea lui Dumnezeu, ce vom spune? Este nedrept Dumnezeu care provoacă mânie? Vorbesc așa cum vorbesc oamenii.

\VS{6}Să nu se întâmple niciodată! Căci atunci cum va judeca Dumnezeu lumea?

\VS{7}Căci dacă adevărul lui Dumnezeu, prin minciuna mea, a abundat spre gloria lui, de ce sunt și eu judecat tot ca păcătos?

\VS{8}De ce nu (așa cum se spune în mod calomnios și cum afirmă unii că spunem): “Să facem răul, ca să vină binele?”? Cei care spun așa sunt condamnați pe bună dreptate.


\par }{\PP \VS{9}Și atunci, ce se întâmplă? Suntem noi mai buni decât ei? Nu, în niciun caz. Căci noi i-am avertizat anterior atât pe iudei, cât și pe greci, că toți sunt sub păcat.

\VS{10}După cum este scris,

\par }{\Q “Nu este nimeni neprihănit;

\par }{\QB nu, nici unul.


\par }{\Q \VS{11}Nu este nimeni care să înțeleagă.

\par }{\QB Nu există nimeni care să-L caute pe Dumnezeu.


\par }{\Q \VS{12}Toți s-au îndepărtat.

\par }{\QB Împreună au devenit neprofitabile.

\par }{\Q Nu există nimeni care să facă binele,

\par }{\QB nu, nici măcar unul.”
\FTNT{†}{{\FR 3:12 }{\FT{{\CHAR{Psalmul 14:1-3}};
{\CHAR{53:1-3}}; Eclesiastul 7:20}}}


\par }{\Q \VS{13}“Gâtul lor este un mormânt deschis.

\par }{\QB Cu limbile lor au folosit înșelăciunea.”
\FTNT{‡}{{\FR 3:13 }{\FT{{\CHAR{Psalmul 5:9}}}}}

\par }{\Q “Otrava viperelor este sub buzele lor.”
\FTNT{§}{{\FR 3:13 }{\FT{{\CHAR{Psalmul 140:3}}}}}


\par }{\QB \VS{14}“Gura lor este plină de blestem și de amărăciune.”
\FTNT{*}{{\FR 3:14 }{\FT{{\CHAR{Psalmul 10:7}}}}}


\par }{\Q \VS{15}Picioarele lor sunt iuți la vărsarea sângelui.


\par }{\QB \VS{16}Distrugerea și nenorocirea sunt în căile lor.


\par }{\QB \VS{17}Calea păcii, ei nu au cunoscut-o.”
\FTNT{†}{{\FR 3:17 }{\FT{{\CHAR{Isaia 59:7-8}}}}}


\par }{\Q \VS{18}“Nu este frică de Dumnezeu înaintea ochilor lor.”
\FTNT{‡}{{\FR 3:18 }{\FT{{\CHAR{Psalmul 36:1}}}}}


\par }{\PP \VS{19}Și știm că tot ce spune Legea vorbește celor ce sunt sub Lege, pentru ca orice gură să fie închisă și toată lumea să fie supusă judecății lui Dumnezeu.

\VS{20}Fiindcă prin faptele Legii, nicio făptură nu va fi îndreptățită înaintea Lui; căci prin Lege vine cunoașterea păcatului.


\par }{\PP \VS{21}Dar acum, afară de Lege, s-a descoperit o neprihănire a lui Dumnezeu, mărturisită prin Lege și prin prooroci,

\VS{22}neprihănirea lui Dumnezeu prin credința în Isus Hristos, pentru toți și pentru toți cei ce cred. Căci nu este nicio deosebire,

\VS{23}căci toți au păcătuit și sunt lipsiți de gloria lui Dumnezeu;

\VS{24}fiind îndreptățiți gratuit prin harul său, prin răscumpărarea care este în Isus Cristos,

\VS{25}pe care Dumnezeu l-a trimis ca jertfă\FTNT{§}{{\FR 3:25 }{\FT{sau, o ispășire}}} de ispășire prin credința în sângele său, pentru o demonstrație a dreptății sale prin trecerea peste păcatele anterioare, în îngăduința lui Dumnezeu;

\VS{26}pentru a demonstra dreptatea sa în timpul prezent, pentru ca el însuși să fie drept și justițiarul celui care are credință în Isus.


\par }{\PP \VS{27}Atunci unde este lăudăroșenia? Ea este exclusă. Prin ce fel de lege? A faptelor? Nu, ci printr-o lege a credinței.

\VS{28}Susținem deci că omul este justificat prin credință, fără faptele legii.

\VS{29}Sau Dumnezeu este oare numai Dumnezeul iudeilor? Nu este el Dumnezeul și al neamurilor? Ba da, și al neamurilor,

\VS{30}de vreme ce, într-adevăr, există un singur Dumnezeu care va justifica pe cei tăiați împrejur prin credință și pe cei netăiați împrejur prin credință.


\par }{\PP \VS{31}În acest caz, anulăm noi Legea prin credință? Să nu se întâmple niciodată! Nu, noi stabilim legea.



\par }\Chap{4}{\PP \VerseOne{1}Atunci ce vom spune că Avraam, strămoșul nostru, a găsit după trup?

\VS{2}Căci, dacă Avraam a fost îndreptățit prin fapte, are cu ce să se laude, dar nu față de Dumnezeu.

\VS{3}Căci ce spune Scriptura? “Avraam a crezut pe Dumnezeu și i s-a socotit ca neprihănire.”
\FTNT{*}{{\FR 4:3 }{\FT{Geneza 15:6}}}

\VS{4}Or, pentru cel care lucrează, răsplata nu este socotită ca un har, ci ca ceva datorat.

\VS{5}Dar celui care nu lucrează, ci crede în Cel care justifică pe cei neevlavioși, credința lui este socotită drept dreptate.

\VS{6}Așa cum și David pronunță binecuvântarea asupra omului căruia Dumnezeu îi socotește dreptate în afară de fapte:


\par }{\Q \VS{7}“Ferice de cei cărora li se iartă fărădelegile,

\par }{\QB ale căror păcate sunt acoperite.


\par }{\Q \VS{8}Binecuvântat este omul pe care Domnul nu-l va acuza de păcat.”
\FTNT{†}{{\FR 4:8 }{\FT{{\CHAR{Psalmul 32:1-2}}}}}


\par }{\PP \VS{9}Este deci binecuvântarea aceasta numai pentru cei tăiați împrejur, sau și pentru cei netăiați împrejur? Căci noi spunem că credința i-a fost socotită lui Avraam ca neprihănire.

\VS{10}Cum a fost deci socotită? Când era în circumcizie, sau în necircumcizie? Nu în circumcizie, ci în necircumcizie.

\VS{11}El a primit semnul circumciziei, ca pecete a neprihănirii credinței pe care o avea când era în necircumcizie, ca să fie tatăl tuturor celor ce cred, chiar dacă sunt în necircumcizie, pentru ca și lor să li se socotească neprihănirea.

\VS{12}El este tatăl circumciziei pentru cei care nu numai că sunt din circumcizie, dar care și umblă pe urmele acelei credințe a tatălui nostru Avraam, pe care el a avut-o în necircumcizie.


\par }{\PP \VS{13}Căci făgăduința făcută lui Avraam și urmașilor lui, că el va fi moștenitorul lumii, nu a fost făcută prin lege, ci prin neprihănirea credinței.

\VS{14}Căci, dacă moștenitorii sunt cei ce țin de Lege, credința se anulează și făgăduința rămâne fără efect.

\VS{15}Căci legea produce mânie; căci unde nu este lege, nu este nici neascultare.


\par }{\PP \VS{16}De aceea este din credință, ca să fie după har, pentru ca făgăduința să fie sigură pentru toți urmașii, nu numai pentru cei din lege, ci și pentru cei din credința lui Avraam, care este tatăl nostru al tuturor.

\VS{17}După cum este scris: “Te-am făcut tatăl multor neamuri”.\FTNT{‡}{{\FR 4:17 }{\FT{Geneza 17:5}}} Aceasta în prezența celui în care a crezut: Dumnezeu, care dă viață celor morți și numește lucrurile care nu sunt, ca și cum ar fi.

\VS{18}Împotriva nădejdii, Avraam a crezut în nădejde, ca să devină tatăl multor neamuri, potrivit cu ceea ce fusese spus: “Așa va fi urmașul tău”.
\FTNT{§}{{\FR 4:18 }{\FT{Geneza 15:5}}}

\VS{19}Fără să slăbească în credință, el nu s-a gândit la propriul său trup, care era deja uzat (el fiind în vârstă de aproape o sută de ani), și la moartea pântecelui Sarei.

\VS{20}Cu toate acestea, privind la promisiunea lui Dumnezeu, el nu s-a clătinat din necredință, ci s-a întărit prin credință, dând slavă lui Dumnezeu

\VS{21}și fiind pe deplin încredințat că ceea ce promisese, era și în stare să împlinească.

\VS{22}De aceea și acest lucru i-a fost “creditat ca neprihănire”.
\FTNT{*}{{\FR 4:22 }{\FT{Geneza 15:6}}}

\VS{23}Or, nu a fost scris că i s-a socotit numai pentru el,

\VS{24}ci și pentru noi, cărora li se va socoti, care credem în cel care l-a înviat din morți pe Isus, Domnul nostru,

\VS{25}care a fost dat pentru greșelile noastre și a înviat pentru îndreptățirea noastră.



\par }\Chap{5}{\PP \VerseOne{1}Deci, fiind îndreptățiți prin credință, avem pace cu Dumnezeu, prin Domnul nostru Isus Hristos,

\VS{2}prin care, de asemenea, avem acces prin credință la acest har în care stăm. Ne bucurăm în speranța gloriei lui Dumnezeu.

\VS{3}Și nu numai atât, ci ne bucurăm și în suferințele noastre, știind că suferința produce perseverență;

\VS{4}iar perseverența, caracter dovedit; iar caracterul dovedit, speranță;

\VS{5}iar speranța nu ne dezamăgește, pentru că dragostea lui Dumnezeu a fost revărsată în inimile noastre prin Duhul Sfânt care ne-a fost dat.


\par }{\PP \VS{6}Căci, pe când eram noi încă slabi, Hristos a murit la vremea potrivită pentru cei nelegiuiți.

\VS{7}Căci cu greu va muri cineva pentru un om neprihănit. Totuși, poate că pentru un om bun cineva chiar va îndrăzni să moară.

\VS{8}Dar Dumnezeu Își recomandă dragostea față de noi, prin faptul că, pe când eram încă păcătoși, Hristos a murit pentru noi.


\par }{\PP \VS{9}Cu atât mai mult, fiind acum îndreptățiți prin sângele Lui, vom fi mântuiți prin El de mânia lui Dumnezeu.

\VS{10}Căci dacă, pe când eram vrăjmași, am fost împăcați cu Dumnezeu prin moartea Fiului Său, cu atât mai mult, fiind împăcați, vom fi mântuiți prin viața lui.


\par }{\PP \VS{11}Și nu numai atât, ci ne și bucurăm în Dumnezeu prin Domnul nostru Isus Hristos, prin care am primit acum împăcarea.

\VS{12}Așadar, după cum păcatul a intrat în lume printr-un singur om și moartea prin păcat, tot așa și moartea a trecut la toți oamenii, pentru că toți au păcătuit.

\VS{13}Căci, până la apariția Legii, păcatul era în lume; dar păcatul nu este acuzat când nu există Legea.

\VS{14}Cu toate acestea, moartea a domnit de la Adam până la Moise, chiar și peste cei ale căror păcate nu au fost ca neascultarea lui Adam, care este o prefigurare a celui ce avea să vină.


\par }{\PP \VS{15}Dar darul gratuit nu este ca și fărădelegea. Căci, dacă prin greșeala unuia singur au murit cei mulți, cu atât mai mult harul lui Dumnezeu și darul prin harul unui singur om, Isus Hristos, a abundat pentru cei mulți.

\VS{16}Darul nu este ca prin unul singur care a păcătuit; căci judecata a venit prin unul singur spre osândă, dar darul gratuit a urmat multor fărădelegi spre îndreptățire.

\VS{17}Căci, dacă prin fărădelegea unuia singur a domnit moartea prin unul singur, cu atât mai mult cei care primesc abundența harului și a darului neprihănirii vor domni în viață prin unul singur, Isus Hristos.


\par }{\PP \VS{18}Astfel, după cum, printr-o singură fărădelege, toți oamenii au fost condamnați, tot așa, printr-o singură faptă de dreptate, toți oamenii au fost îndreptățiți pentru viață.

\VS{19}Căci, după cum prin neascultarea unui singur om, mulți au fost făcuți păcătoși, tot așa, prin ascultarea unuia singur, mulți vor fi făcuți neprihăniți.

\VS{20}Legea a venit pentru ca fărădelegea să abunde; dar acolo unde a abundat păcatul, harul a abundat și mai mult,

\VS{21}pentru ca, după cum păcatul a domnit în moarte, tot așa harul să domnească prin dreptate pentru viața veșnică, prin Isus Hristos, Domnul nostru.



\par }\Chap{6}{\PP \VerseOne{1}Ce să spunem atunci? Să continuăm în păcat, pentru ca harul să abunde?

\VS{2}Să nu se întâmple niciodată! Noi, care am murit față de păcat, cum am mai putea trăi în el?

\VS{3}Sau nu știți că noi toți cei care am fost botezați în Hristos Isus am fost botezați în moartea lui?

\VS{4}Am fost deci îngropați împreună cu el prin botez în moarte, pentru ca, după cum Hristos a înviat din morți prin gloria Tatălui, tot așa și noi să umblăm într-o viață nouă.


\par }{\PP \VS{5}Căci, dacă ne-am unit cu El în asemănarea morții Lui, vom fi și noi părtași la învierea Lui,

\VS{6}știind că omul nostru cel vechi a fost răstignit împreună cu El, pentru ca trupul păcatului să fie desființat, ca să nu mai fim robi ai păcatului.

\VS{7}Căci cel care a murit a fost eliberat de păcat.

\VS{8}Dar dacă am murit împreună cu Hristos, credem că și noi vom trăi împreună cu el,

\VS{9}știind că Hristos, înviind din morți, nu mai moare. Moartea nu mai are stăpânire asupra lui!

\VS{10}Căci moartea în care a murit, a murit o singură dată pentru păcat; dar viața pe care o trăiește, trăiește pentru Dumnezeu.

\VS{11}Astfel, considerați-vă și voi înșivă ca fiind morți față de păcat, dar vii pentru Dumnezeu în Hristos Isus, Domnul nostru.


\par }{\PP \VS{12}De aceea, nu lăsați păcatul să domnească în trupul vostru muritor, ca să-l ascultați în poftele lui.

\VS{13}De asemenea, nu vă prezentați membrele voastre păcatului ca instrumente ale nedreptății, ci prezentați-vă lui Dumnezeu ca vii din morți, iar membrele voastre ca instrumente ale dreptății pentru Dumnezeu.

\VS{14}Căci păcatul nu va avea stăpânire asupra voastră, pentru că nu sunteți sub lege, ci sub har.


\par }{\PP \VS{15}Și atunci, ce se întâmplă? Să păcătuim pentru că nu suntem sub lege, ci sub har? Să nu se întâmple niciodată!

\VS{16}Nu știți că, atunci când vă prezentați ca niște robi și ascultați de cineva, sunteți robii celui de care ascultați, fie de păcat spre moarte, fie de ascultare spre dreptate?

\VS{17}Dar mulțumim lui Dumnezeu că, în timp ce erați robi ai păcatului, ați devenit ascultători din inimă față de acea formă de învățătură la care ați fost predați.

\VS{18}Fiind eliberați de păcat, ați devenit robi ai neprihănirii.


\par }{\PP \VS{19}Vorbesc în termeni omenești, din pricina slăbiciunii cărnii voastre; căci, după cum v-ați prezentat mădularele ca robi ai necurăției și ai răutății peste răutăți, tot așa acum prezentați mădularele voastre ca robi ai neprihănirii pentru sfințire.

\VS{20}Căci atunci când erați robi ai păcatului, erați liberi de neprihănire.

\VS{21}Ce roade aveați deci atunci în lucrurile de care vă rușinați acum? Căci sfârșitul acelor lucruri este moartea.

\VS{22}Dar acum, fiind eliberați de păcat și deveniți robi ai lui Dumnezeu, aveți rodul sfințirii și rezultatul vieții veșnice.

\VS{23}Căci plata păcatului este moartea, dar darul gratuit al lui Dumnezeu este viața veșnică în Hristos Isus, Domnul nostru.



\par }\Chap{7}{\PP \VerseOne{1}Sau nu știți voi, fraților, căci vorbesc cu oameni care cunosc Legea, că Legea stăpânește asupra omului cât trăiește el?

\VS{2}Căci femeia care are un soț este legată prin lege de soț cât trăiește el, dar dacă soțul moare, ea este eliberată de legea soțului.

\VS{3}Deci, dacă, cât timp trăiește soțul, ea se unește cu un alt bărbat, ar fi numită adulteră. Dar dacă soțul moare, ea este eliberată de legea soțului, astfel că nu este adulteră, deși este unită cu un alt bărbat.

\VS{4}Așadar, frații mei, și voi ați fost făcuți morți față de Lege prin trupul lui Hristos, ca să vă uniți cu altul, cu Cel înviat din morți, ca să producem roade pentru Dumnezeu.

\VS{5}Căci, atunci când eram în carne, patimile păcătoase care erau prin lege lucrau în membrele noastre pentru a produce roade spre moarte.

\VS{6}Dar acum am fost eliberați de Lege, după ce am murit față de cea în care eram ținuți, astfel încât slujim în noutatea duhului, și nu în vechimea literei.


\par }{\PP \VS{7}Ce vom spune atunci? Este legea păcat? Să nu fie niciodată! Cu toate acestea, nu aș fi cunoscut păcatul decât prin lege. Căci nu aș fi cunoscut pofta dacă legea nu ar fi spus: “Să nu poftești”.
\FTNT{*}{{\FR 7:7 }{\FT{Exodul 20:17; Deuteronomul 5:21}}}

\VS{8}Dar păcatul, găsind prilej prin poruncă, a produs în mine tot felul de pofte. Căci, în afara legii, păcatul este mort.

\VS{9}În afară de Lege, am fost odată viu, dar când a venit porunca, păcatul a reînviat și am murit.

\VS{10}Porunca, care era pentru viață, am aflat că este pentru moarte;

\VS{11}căci păcatul, găsind prilej prin poruncă, m-a înșelat și prin ea m-a omorât.

\VS{12}De aceea, Legea este cu adevărat sfântă și porunca sfântă, dreaptă și bună.


\par }{\PP \VS{13}Deci ceea ce este bun a devenit pentru mine moarte? Să nu fie niciodată! Dar păcatul, ca să se arate că este păcat, producea moarte în mine prin ceea ce este bun, pentru ca prin poruncă păcatul să devină peste măsură de păcătos.

\VS{14}Căci știm că Legea este spirituală, dar eu sunt trupesc, vândut sub păcat.

\VS{15}Căci nu înțeleg ce fac. Căci nu practic ceea ce doresc să fac, ci ceea ce urăsc, aceea fac.

\VS{16}Dar dacă ceea ce nu doresc, aceea fac, consimt Legii că este bună.

\VS{17}Deci acum nu mai sunt eu cel care o face, ci păcatul care locuiește în mine.

\VS{18}Căci știu că în mine, adică în carnea mea, nu locuiește nimic bun. Căci dorința este prezentă la mine, dar nu o găsesc făcând ceea ce este bun.

\VS{19}Căci binele pe care-l doresc, nu-l fac; dar răul pe care nu-l doresc, pe acela îl practic.

\VS{20}Dar dacă ceea ce nu doresc, fac, nu mai sunt eu cel care face, ci păcatul care locuiește în mine.

\VS{21}Găsesc deci legea că, în timp ce eu doresc să fac binele, răul este prezent.

\VS{22}Căci eu mă bucur de legea lui Dumnezeu după cele lăuntrice,

\VS{23}dar văd în mădularele mele o lege diferită, care se luptă împotriva legii minții mele și mă aduce în captivitate sub legea păcatului care este în mădularele mele.

\VS{24}Ce om nenorocit sunt eu! Cine mă va izbăvi din trupul acestei morți?

\VS{25}Mulțumesc lui Dumnezeu prin Isus Hristos, Domnul nostru! Așadar, cu mintea, eu însumi slujesc legii lui Dumnezeu, dar cu carnea, legii păcatului.



\par }\Chap{8}{\PP \VerseOne{1}Așadar, acum nu este nici o condamnare pentru cei ce sunt în Hristos Isus, care nu umblă după trup, ci după Duhul.
\FTNT{*}{{\FR 8:1 }{\FT{TR omite “Dumnezeu”}}}

\VS{2}Căci legea Duhului de viață în Hristos Isus m-a făcut liber de legea păcatului și a morții.

\VS{3}Căci ceea ce nu putea face legea, fiindcă era slabă prin carne, Dumnezeu a făcut, trimițându-și propriul Fiu în chip de carne păcătoasă și pentru păcat, a condamnat păcatul în carne,

\VS{4}pentru ca rânduiala legii să fie împlinită în noi, care nu umblăm după carne, ci după Duhul.

\VS{5}Căci cei care trăiesc după trup își pun mintea la lucrurile trupului, dar cei care trăiesc după Duhul, la lucrurile Duhului.

\VS{6}Căci mintea cărnii este moarte, dar mintea Duhului este viață și pace;

\VS{7}pentru că mintea cărnii este ostilă lui Dumnezeu, deoarece nu se supune legii lui Dumnezeu și nici nu poate fi.

\VS{8}Cei care sunt în carne nu pot să placă lui Dumnezeu.


\par }{\PP \VS{9}Dar voi nu sunteți în carne, ci în Duh, dacă Duhul lui Dumnezeu locuiește în voi. Dar dacă cineva nu are Duhul lui Hristos, nu este al lui.

\VS{10}Dacă Hristos este în voi, trupul este mort din pricina păcatului, dar spiritul este viu din pricina neprihănirii.

\VS{11}Dar dacă Duhul celui care l-a înviat pe Isus din morți locuiește în voi, cel care l-a înviat pe Isus Hristos din morți va da viață și trupurilor voastre muritoare, prin Duhul Său care locuiește în voi.


\par }{\PP \VS{12}Deci, fraților, nu suntem datori să trăim după trup, ca să trăim după trup.

\VS{13}Căci, dacă trăiți după trup, trebuie să muriți; dar dacă, prin Duhul, faceți să moară faptele trupului, veți trăi.

\VS{14}Căci toți cei ce sunt călăuziți de Duhul lui Dumnezeu, aceștia sunt copiii lui Dumnezeu.

\VS{15}Căci nu ați primit din nou duhul robiei pentru frică, ci ați primit Duhul adopției, prin care strigăm: “Abba! Tată!”


\par }{\PP \VS{16}Însuși Duhul mărturisește cu duhul nostru că suntem copii ai lui Dumnezeu;

\VS{17}și, dacă suntem copii, suntem și moștenitori, moștenitori ai lui Dumnezeu și împreună-moștenitori cu Hristos, dacă suferim împreună cu El, ca să fim și noi proslăviți împreună cu El.


\par }{\PP \VS{18}Căci consider că suferințele de acum nu sunt vrednice de comparație cu slava care se va arăta față de noi.

\VS{19}Căci creația așteaptă cu nerăbdare arătarea copiilor lui Dumnezeu.

\VS{20}Căci creația a fost supusă deșertăciunii, nu din voia ei, ci din cauza celui care a supus-o, în speranța

\VS{21}că și creația însăși va fi eliberată din robia stricăciunii în libertatea gloriei copiilor lui Dumnezeu.

\VS{22}Căci știm că întreaga creație geme și se chinuiește în dureri împreună până acum.

\VS{23}Nu numai noi înșine, ci și noi înșine, care avem primele roade ale Duhului, chiar și noi înșine gemem în noi înșine, așteptând adopția, răscumpărarea trupului nostru.

\VS{24}Căci am fost mântuiți în speranță, dar speranța care se vede nu este speranță. Căci cine speră în ceea ce vede?

\VS{25}Dar dacă sperăm în ceea ce nu vedem, așteptăm cu răbdare.


\par }{\PP \VS{26}Tot așa și Duhul Sfânt ne ajută în slăbiciunile noastre, pentru că nu știm să ne rugăm cum trebuie. Dar însuși Duhul mijlocește pentru noi cu gemete care nu pot fi rostite.

\VS{27}Cel care cercetează inimile știe ce este în mintea Duhului, pentru că el mijlocește pentru sfinți, potrivit lui Dumnezeu.


\par }{\PP \VS{28}Știm că toate lucrurile lucrează împreună spre binele celor ce iubesc pe Dumnezeu, al celor ce sunt chemați după planul Lui.

\VS{29}Căci pe cei pe care i-a cunoscut mai dinainte, i-a și predestinat să fie conformați cu chipul Fiului Său, pentru ca El să fie cel dintâi născut între mulți frați.

\VS{30}Pe cei pe care i-a predestinat, pe aceia i-a și chemat. Pe cei pe care i-a chemat, pe aceia i-a și îndreptățit. Pe cei pe care i-a justificat, pe aceia i-a și glorificat.


\par }{\PP \VS{31}Ce să spunem despre aceste lucruri? Dacă Dumnezeu este pentru noi, cine poate fi împotriva noastră?

\VS{32}El, care nu și-a cruțat propriul Fiu, ci L-a predat pentru noi toți, cum nu ne-ar da și El, împreună cu El, toate lucrurile?

\VS{33}Cine ar putea aduce o acuzație împotriva aleșilor lui Dumnezeu? Dumnezeu este cel care justifică.

\VS{34}Cine este cel care condamnă? Este Hristos, care a murit, ba mai mult, care a înviat din morți, care este la dreapta lui Dumnezeu, care și el mijlocește pentru noi.


\par }{\PP \VS{35}Cine ne va despărți pe noi de dragostea lui Hristos? Ar putea oare opresiunea, sau chinul, sau persecuția, sau foametea, sau goliciunea, sau primejdia, sau sabia?

\VS{36}Așa cum este scris,

\par }{\Q “De dragul tău suntem uciși toată ziua.

\par }{\QB Am fost socotiți ca niște oi pentru măcelărie.”


\par }{\MM \VS{37}Nu, în toate aceste lucruri suntem mai mult decât biruitori prin Cel ce ne-a iubit.

\VS{38}Căci sunt încredințat că nici moartea, nici viața, nici îngerii, nici stăpânirile, nici lucrurile prezente, nici cele viitoare, nici puterile,

\VS{39}nici înălțimea, nici adâncimea, nici vreun alt lucru creat nu va putea să ne despartă de dragostea lui Dumnezeu, care este în Hristos Isus, Domnul nostru.



\par }\Chap{9}{\PP \VerseOne{1}Eu spun adevărul în Hristos. Nu mint, conștiința mea mărturisește cu mine în Duhul Sfânt

\VS{2}că am o mare întristare și o durere neîncetată în inima mea.

\VS{3}Căci aș putea să doresc să fiu eu însumi blestemat de Cristos, din cauza fraților mei, rudele mele după trup

\VS{4}care sunt israeliți; a căror adopție, glorie, legăminte, dăruire a legii, slujire și promisiuni

\VS{5}din care sunt părinții și din care provine Cristos în ceea ce privește trupul, care este peste toate, Dumnezeu, binecuvântat în veci. Amin.


\par }{\PP \VS{6}Dar nu este ca și cum Cuvântul lui Dumnezeu ar fi fost zădărnicit. Căci nu toți cei din Israel sunt din Israel.

\VS{7}și nici, pentru că sunt urmașii lui Avraam, nu sunt toți copii. Ci, “urmașii voștri vor fi socotiți ca de la Isaac”.

\VS{8}Adică, nu copiii cărnii sunt copiii lui Dumnezeu, ci copiii promisiunii sunt socotiți moștenitori.

\VS{9}Căci acesta este un cuvânt al promisiunii: “La timpul hotărât voi veni și Sara va avea un fiu”.

\VS{10}Și nu numai atât, ci și Rebeca a conceput un fiu, pe tatăl nostru Isaac.

\VS{11}Căci, nefiind încă născută, și neavând nimic bun sau rău, pentru ca scopul lui Dumnezeu, potrivit alegerii, să rămână în picioare, nu din fapte, ci de la Cel care cheamă,

\VS{12}i s-a spus: “Cel mai mare va sluji celui mai mic”.

\VS{13}Așa cum este scris: “Pe Iacov l-am iubit, dar pe Esau l-am urât”.


\par }{\PP \VS{14}Ce vom spune atunci? Există nedreptate la Dumnezeu? Să nu fie niciodată!

\VS{15}Căci El a zis lui Moise: “Voi avea milă de cine am milă și voi avea milă de cine am milă”.

\VS{16}Deci, nu este a celui care vrea, nici a celui care aleargă, ci a lui Dumnezeu care are milă.

\VS{17}Căci Scriptura îi spune lui Faraon: “Tocmai pentru aceasta am făcut să te înalț, ca să-mi arăt în tine puterea Mea și ca numele Meu să fie vestit pe tot pământul.”

\VS{18}Așadar, El are milă de cine vrea și împietrește pe cine vrea.


\par }{\PP \VS{19}Atunci îmi veți zice: “De ce mai găsește el greșeală? Căci cine se împotrivește voinței lui?”

\VS{20}Dar, într-adevăr, omule, cine ești tu ca să răspunzi împotriva lui Dumnezeu? Oare lucrul format îl va întreba pe cel care l-a format: “De ce m-ai făcut așa?”

\VS{21}Sau nu are olarul un drept asupra lutului, din aceeași bucată să facă dintr-o parte un vas pentru cinste și din alta pentru dezonoare?

\VS{22}Și dacă Dumnezeu, vrând să-și arate mânia și să-și facă cunoscută puterea, a îndurat cu multă răbdare vasele mâniei pregătite pentru distrugere,

\VS{23}și ca să facă cunoscută bogăția gloriei sale pe vasele milei, pe care le-a pregătit mai dinainte pentru glorie—

\VS{24}pe noi, pe care ne-a chemat și pe noi, nu numai dintre iudei, ci și dintre neamuri?

\VS{25}După cum spune și în Osea,

\par }{\Q “Îi voi numi “poporul meu”, pe cei care nu erau poporul meu;

\par }{\QB și “iubitul” ei, care nu era iubit.”


\par }{\Q \VS{26}În locul în care li s-a spus: “Voi nu sunteți poporul Meu”.

\par }{\QB acolo vor fi numiți “copiii Dumnezeului celui viu”.”


\par }{\PP \VS{27}Isaia strigă cu privire la Israel,

\par }{\Q “Dacă numărul copiilor lui Israel este ca nisipul mării,

\par }{\QB rămășița este cea care va fi salvată;


\par }{\Q \VS{28}Căci El va isprăvi lucrarea și o va curma cu dreptate,

\par }{\QB pentru că Domnul va face o lucrare scurtă pe pământ.”


\par }{\PP \VS{29}După cum a spus Isaia mai înainte,

\par }{\Q “Dacă Domnul oștirilor nu ne-ar fi lăsat o sămânță,

\par }{\QB am fi devenit ca Sodoma,

\par }{\QB și ar fi fost făcut ca Gomora.”


\par }{\PP \VS{30}Ce vom spune atunci? Că neamurile, care nu urmăreau neprihănirea, au ajuns la neprihănire, la neprihănirea care vine din credință;

\VS{31}dar Israel, care urmărea o lege a neprihănirii, nu a ajuns la legea neprihănirii.

\VS{32}De ce? Pentru că nu au căutat-o prin credință, ci ca și cum ar fi fost prin faptele legii. S-au împiedicat de piatra de poticnire,

\VS{33}așa cum este scris,

\par }{\Q “Iată, pun în Sion o piatră de poticnire și o stâncă de jignire;

\par }{\QB și nimeni care crede în El nu va fi dezamăgit.”



\par }\Chap{10}{\PP \VerseOne{1}Fraților, dorința inimii mele și rugăciunea mea către Dumnezeu este pentru Israel, ca să fie mântuit.

\VS{2}Căci mărturisesc despre ei că au un zel pentru Dumnezeu, dar nu potrivit cu cunoștința.

\VS{3}Pentru că, ignorând dreptatea lui Dumnezeu și căutând să își stabilească propria dreptate, nu s-au supus dreptății lui Dumnezeu.

\VS{4}Căci Hristos este împlinirea Legii, ca neprihănire pentru oricine crede.


\par }{\PP \VS{5}Căci Moise scrie despre dreptatea Legii: “Cel ce le împlinește va trăi prin ele.”

\VS{6}Dar neprihănirea care vine din credință spune următoarele: “Nu spuneți în inima voastră: “Cine se va sui în cer?”. (adică să-l coboare pe Cristos);

\VS{7}sau: “Cine se va coborî în abis?” (adică să-l coboare pe Cristos); sau: “Cine se va coborî în abis?”. (adică să-l aducă pe Hristos din morți).”

\VS{8}Dar ce spune? “Cuvântul este aproape de tine, în gura ta și în inima ta”, adică cuvântul credinței pe care îl predicăm:

\VS{9}că, dacă vei mărturisi cu gura ta că Isus este Domnul și dacă vei crede în inima ta că Dumnezeu l-a înviat din morți, vei fi mântuit.

\VS{10}Căci cu inima se crede, ceea ce duce la neprihănire, iar cu gura se face mărturisirea, ceea ce duce la mântuire.

\VS{11}Căci Scriptura spune: “Oricine crede în el nu va fi dezamăgit”.


\par }{\PP \VS{12}Căci nu este nici o deosebire între iudei și greci, căci același Domn este Domnul tuturor și este bogat pentru toți cei ce-L cheamă.

\VS{13}Căci: “Oricine va invoca numele Domnului va fi mântuit”.

\VS{14}Cum îl vor invoca, așadar, pe cel în care nu au crezut? Cum vor crede în cel pe care nu l-au auzit? Cum vor auzi fără un predicator?

\VS{15}Și cum vor propovădui ei dacă nu vor fi trimiși? După cum este scris:

\par }{\Q “Cât de frumoase sunt picioarele celor care propovăduiesc vestea bună a păcii,

\par }{\QB care aduceți vești bune!”


\par }{\PP \VS{16}Dar nu toți au ascultat vestea cea bună. Căci Isaia spune: “Doamne, cine a crezut în vestea noastră?”.

\VS{17}Așadar, credința vine prin auz, iar auzul prin cuvântul lui Dumnezeu.

\VS{18}Dar eu zic: nu au auzit ei? Ba da, cu siguranță că da,

\par }{\Q “Sunetul lor s-a răspândit pe tot pământul,

\par }{\QB cuvintele lor până la marginile lumii.”


\par }{\PP \VS{19}Dar eu întreb: nu știa Israel? Mai întâi Moise spune,

\par }{\Q “Vă voi provoca la gelozie cu ceea ce nu este neam.

\par }{\QB Te voi înfuria cu un neam lipsit de înțelegere.”


\par }{\PP \VS{20}Isaia este foarte îndrăzneț și spune,

\par }{\Q “Am fost găsit de cei care nu m-au căutat.

\par }{\QB Am fost dezvăluit celor care nu m-au cerut.”


\par }{\PP \VS{21}Dar despre Israel spune: “Toată ziua mi-am întins mâinile spre un popor neascultător și răzvrătit.”



\par }\Chap{11}{\PP \VerseOne{1}Întreb atunci: A respins Dumnezeu pe poporul Său? Fie ca niciodată să nu fie așa! Căci și eu sunt israelit, descendent al lui Avraam, din seminția lui Beniamin.

\VS{2}Dumnezeu nu și-a respins poporul său, pe care l-a cunoscut dinainte. Sau nu știți ce spune Scriptura despre Ilie? Cum îl imploră el pe Dumnezeu împotriva lui Israel:

\VS{3}“Doamne, ei au ucis pe profeții Tăi. Au dărâmat altarele Tale. Eu am rămas singur, iar ei îmi caută viața”.

\VS{4}Dar cum îi răspunde Dumnezeu? “Mi-am rezervat șapte mii de oameni care nu și-au plecat genunchiul în fața lui Baal.”

\VS{5}Tot așa, și în acest timp prezent există o rămășiță, potrivit alegerii harului.

\VS{6}Și dacă este prin har, atunci nu mai este prin fapte; altfel, harul nu mai este har. Dar dacă este din fapte, nu mai este har; altfel fapta nu mai este faptă.


\par }{\PP \VS{7}Și atunci, ce se întâmplă? Ceea ce caută Israel, nu a obținut, ci au obținut-o cei aleși, iar ceilalți s-au împietrit.

\VS{8}După cum este scris: “Dumnezeu le-a dat un duh de stupoare, ochi ca să nu vadă și urechi ca să nu audă, până în ziua de azi”.


\par }{\PP \VS{9}David spune,

\par }{\Q “Masa lor să se facă o cursă, o capcană,

\par }{\QB o piatră de poticnire și o pedeapsă pentru ei.


\par }{\Q \VS{10}Să li se întunece ochii, ca să nu vadă.

\par }{\QB Întotdeauna țineți-le spatele îndoit.”


\par }{\PP \VS{11}Întreb, deci, dacă s-au împiedicat ca să cadă? Să nu se întâmple niciodată! Dar, prin căderea lor, mântuirea a ajuns la neamuri, ca să le provoace gelozia.

\VS{12}Or, dacă căderea lor este bogăția lumii și pierderea lor este bogăția neamurilor, cu cât mai mult plinătatea lor!


\par }{\PP \VS{13}Căci Eu vă vorbesc vouă, care sunteți neamuri. De vreme ce, fiind apostol al neamurilor, îmi slăvesc slujba,

\VS{14}dacă, prin orice mijloc, pot să provoc la gelozie pe cei din carnea mea și să salvez pe unii dintre ei.

\VS{15}Căci, dacă respingerea lor este împăcarea lumii, ce ar fi acceptarea lor, dacă nu viața din morți?


\par }{\PP \VS{16}Dacă cel dintâi fruct este sfânt, la fel este și masa. Dacă rădăcina este sfântă, și ramurile sunt sfinte.

\VS{17}Dar dacă unele dintre ramuri au fost rupte, iar tu, fiind un măslin sălbatic, ai fost altoit printre ele și ai devenit părtaș cu ele la rădăcina și la bogăția măslinului,

\VS{18}nu te lăuda cu ramurile. Dar dacă te lauzi, adu-ți aminte că nu tu susții rădăcina, ci rădăcina te susține pe tine.

\VS{19}Veți spune atunci: “Ramurile au fost rupte, ca să fiu altoit eu”.

\VS{20}Adevărat; prin necredința lor au fost rupte, iar voi stați în picioare prin credința voastră. Nu fiți îngâmfați, ci temători;

\VS{21}căci dacă Dumnezeu nu a cruțat ramurile naturale, nici pe voi nu vă va cruța.

\VS{22}Vedeți, așadar, bunătatea și severitatea lui Dumnezeu. Față de cei care au căzut, severitate; dar față de voi, bunătate, dacă rămâneți în bunătatea lui; altfel, și voi veți fi stârpiți.

\VS{23}Și ei, dacă nu vor continua în necredința lor, vor fi altoiți, căci Dumnezeu este în stare să-i altoiască din nou.

\VS{24}Căci dacă voi ați fost tăiați din ceea ce este prin natură un măslin sălbatic și ați fost altoiți, contrar naturii, într-un măslin bun, cu cât mai mult aceștia, care sunt ramurile naturale, vor fi altoiți în propriul lor măslin?


\par }{\PP \VS{25}Căci nu voiesc, fraților, să nu cunoașteți taina aceasta, ca să nu fiți înțelepți în închipuirile voastre: că Israel a suferit o împietrire parțială, până ce va veni plinătatea neamurilor,

\VS{26}și tot Israelul va fi mântuit. Așa cum este scris,

\par }{\Q “Din Sion va ieși Eliberatorul,

\par }{\QB și va îndepărta nelegiuirea din Iacov.


\par }{\Q \VS{27}Acesta este legământul Meu cu ei,

\par }{\QB când le voi șterge păcatele.”


\par }{\PP \VS{28}În ceea ce privește Buna Vestire, ei sunt vrăjmași din pricina voastră. Dar în ceea ce privește alegerea, ei sunt iubiți din cauza părinților.

\VS{29}Căci darurile și chemarea lui Dumnezeu sunt irevocabile.

\VS{30}Căci, după cum voi, în trecut, ați fost neascultători față de Dumnezeu, dar acum ați obținut îndurare prin neascultarea lor,

\VS{31}tot așa și aceștia au fost neascultători, pentru ca, prin îndurarea arătată vouă, să obțină și ei îndurare.

\VS{32}Căci Dumnezeu i-a legat pe toți de neascultare, ca să aibă milă de toți.


\par }{\PP \VS{33}O, adâncul bogăției înțelepciunii și a cunoștinței lui Dumnezeu! Cât de nepătrunse sunt judecățile Lui și cât de neîntrecute sunt căile Lui, care nu pot fi trasate!


\par }{\Q \VS{34}Căci cine a cunoscut gândul Domnului?

\par }{\QB Sau cine a fost consilierul lui?”


\par }{\Q \VS{35}“Sau cine i-a dat mai întâi,

\par }{\QB și i se va restitui din nou?”


\par }{\PP \VS{36}Căci din El, prin El și pentru El sunt toate lucrurile. A lui să fie gloria în veci! Amin.



\par }\Chap{12}{\PP \VerseOne{1}De aceea, fraților, vă îndemn, prin mila lui Dumnezeu, să vă aduceți trupurile voastre ca o jertfă vie, sfântă, plăcută lui Dumnezeu, care este slujba voastră spirituală.

\VS{2}Nu vă potriviți cu lumea aceasta, ci transformați-vă prin înnoirea minții voastre, ca să puteți dovedi care este voia lui Dumnezeu, bună, plăcută și desăvârșită.


\par }{\PP \VS{3}Pentru că, prin harul care mi-a fost dat, spun fiecăruia dintre voi să nu se socoată mai mult decât trebuie să se socoată, ci să se socoată drept, după cum Dumnezeu a dat fiecăruia o măsură de credință.

\VS{4}Căci, după cum într-un singur trup avem mai multe mădulare și nu toate mădularele au aceeași funcție,

\VS{5}tot așa și noi, care suntem mulți, suntem un singur trup în Hristos și, în mod individual, mădulare unul altuia,

\VS{6}având daruri diferite, potrivit cu harul care ne-a fost dat: dacă profețește, să profețească după măsura credinței noastre;

\VS{7}sau dacă slujește, să ne dăruim slujirii; sau cine învață, să se dăruiască învățăturii;

\VS{8}sau cine îndeamnă, să se dăruiască îndemnului; cine dăruiește, să o facă cu generozitate; cine conduce, cu sârguință; cine face milostenie, cu veselie.


\par }{\PP \VS{9}Dragostea să fie fără ipocrizie. Detestați ceea ce este rău. Să se agațe de ceea ce este bun.

\VS{10}În dragostea fraților, fiți tandri unii cu alții; în cinste, preferați-vă unii pe alții,

\VS{11}nu rămâneți în urmă cu sârguința, fiți fervenți în duh, slujiți Domnului,

\VS{12}bucurați-vă în speranță, rezistați în necazuri, continuați cu stăruință în rugăciune,

\VS{13}contribuiți la nevoile sfinților și fiți darnici în ospitalitate.


\par }{\PP \VS{14}Binecuvântați pe cei ce vă prigonesc; binecuvântați și nu blestemați.

\VS{15}Bucurați-vă cu cei ce se bucură. Plângeți cu cei ce plâng.

\VS{16}Fiți de aceeași părere unii față de alții. Nu vă gândiți la lucruri înalte, ci asociați-vă cu cei umili. Nu fiți înțelepți în închipuirile voastre.

\VS{17}Nu răsplătiți nimănui rău pentru rău. Respectați ceea ce este onorabil în ochii tuturor oamenilor.

\VS{18}Dacă este posibil, atât cât depinde de tine, fii în pace cu toți oamenii.

\VS{19}Nu căutați voi înșivă să vă răzbunați, iubiților, ci lăsați loc mâniei lui Dumnezeu. Căci este scris: “Mie îmi aparține răzbunarea; Eu voi răsplăti, zice Domnul.”

\VS{20}De aceea

\par }{\Q “Dacă dușmanul tău este flămând, dă-i de mâncare.

\par }{\QB Dacă îi este sete, dați-i să bea;

\par }{\QB pentru că, făcând așa, îi vei arunca cărbuni de foc pe cap.”


\par }{\PP \VS{21}Nu vă lăsați biruiți de rău, ci biruiți răul cu binele.



\par }\Chap{13}{\PP \VerseOne{1}Fiecare suflet să se supună autorităților superioare, căci nu există autoritate decât de la Dumnezeu, iar cele ce există sunt rânduite de Dumnezeu.

\VS{2}De aceea, cel care se împotrivește autorității se împotrivește rânduielii lui Dumnezeu; iar cei care se împotrivesc vor primi pentru ei înșiși judecata.

\VS{3}Căci conducătorii nu sunt o teroare pentru fapta bună, ci pentru cea rea. Doriți să nu vă temeți de autoritate? Fă ceea ce este bun, și vei avea parte de laudă din partea autorității,

\VS{4}căci ea este un slujitor al lui Dumnezeu pentru tine, spre bine. Dar dacă faci ceea ce este rău, teme-te, căci nu degeaba poartă sabia, căci este un slujitor al lui Dumnezeu, răzbunător pentru mânia celui care face răul.

\VS{5}De aceea trebuie să vă supuneți, nu numai din cauza mâniei, ci și din cauza conștiinței.

\VS{6}De aceea și voi plătiți impozite, pentru că sunt slujitori ai lui Dumnezeu, care fac în permanență tocmai acest lucru.

\VS{7}De aceea, dați fiecăruia ceea ce îi este dator: dacă datorați impozite, plătiți impozitele; dacă aveți vameși, atunci vameși; dacă aveți respect, atunci respect; dacă aveți onoare, atunci onoare.


\par }{\PP \VS{8}Să nu datorați nimănui nimic, decât să vă iubiți unii pe alții; căci cine iubește pe aproapele său a împlinit Legea.

\VS{9}Căci poruncile: “Să nu comiți adulter”, “Să nu ucizi”, “Să nu furi”, “Să nu poftești” și orice alte porunci ar mai fi, toate se rezumă la acest cuvânt, și anume: “Să iubești pe aproapele tău ca pe tine însuți”.

\VS{10}Dragostea nu face rău aproapelui. Prin urmare, iubirea este împlinirea legii.


\par }{\PP \VS{11}Fă aceasta, știind că a sosit deja vremea să te trezești din somn, căci mântuirea este mai aproape de noi decât atunci când am crezut.

\VS{12}Noaptea este departe, și ziua este aproape. Să ne lepădăm deci de faptele întunericului și să ne îmbrăcăm cu armura luminii.

\VS{13}Să umblăm cum se cuvine, ca în timpul zilei; nu în desfrâu și beție, nu în promiscuitate sexuală și acte de desfrânare, nu în certuri și gelozie.

\VS{14}Ci îmbrăcați-vă în Domnul Isus Hristos și nu faceți niciun fel de provizii pentru carne, pentru poftele ei.



\par }\Chap{14}{\PP \VerseOne{1}Acceptați acum pe cel slab în credință, dar nu pentru certuri de păreri.

\VS{2}Un om are credință să mănânce de toate, dar cel care este slab mănâncă numai legume.

\VS{3}Să nu-l disprețuiască cel care mănâncă pe cel care nu mănâncă. Cel care nu mănâncă să nu-l judece pe cel care mănâncă, pentru că Dumnezeu l-a acceptat.

\VS{4}Cine sunteți voi care judecați pe robul altuia? În fața propriului său stăpân stă sau cade. Da, el va fi făcut să stea în picioare, căci Dumnezeu are puterea de a-l face să stea în picioare.


\par }{\PP \VS{5}Un om consideră că o zi este mai importantă. Altul apreciază fiecare zi la fel. Fiecare să fie pe deplin sigur în mintea sa.

\VS{6}Cine prăznuiește ziua, o prăznuiește pentru Domnul, iar cine nu prăznuiește ziua, nu o prăznuiește pentru Domnul. Cel care mănâncă, mănâncă pentru Domnul, căci mulțumește lui Dumnezeu. Cel care nu mănâncă, pentru Domnul nu mănâncă, și îi mulțumește lui Dumnezeu.

\VS{7}Căci niciunul dintre noi nu trăiește pentru sine și niciunul nu moare pentru sine.

\VS{8}Căci dacă trăim, trăim pentru Domnul. Sau dacă murim, murim pentru Domnul. Așadar, dacă trăim sau murim, suntem ai Domnului.

\VS{9}Căci pentru aceasta Hristos a murit, a înviat și a trăit din nou, ca să fie Domnul celor morți și al celor vii.


\par }{\PP \VS{10}Dar tu, de ce judeci pe fratele tău? Sau tu, de ce îl disprețuiești pe fratele tău? Căci toți vom sta în fața scaunului de judecată al lui Hristos.

\VS{11}Căci este scris:,

\par }{\Q “Pe viața Mea, zice Domnul, orice genunchi se va pleca înaintea Mea.

\par }{\QB Orice limbă va mărturisi lui Dumnezeu”.”


\par }{\PP \VS{12}Astfel, fiecare dintre noi va da socoteală de sine lui Dumnezeu.


\par }{\PP \VS{13}De aceea să nu ne mai judecăm unii pe alții, ci să judecăm mai degrabă aceasta: să nu pună nimeni piedică fratelui său și nici prilej de cădere.

\VS{14}Știu și sunt încredințat în Domnul Isus că nimic nu este necurat de la sine; numai că pentru cel ce consideră ceva necurat, pentru el este necurat.

\VS{15}Totuși, dacă din pricina mâncării fratele tău este mâhnit, nu mai umblați în dragoste. Nu distruge cu mâncarea ta pe cel pentru care a murit Hristos.

\VS{16}Atunci nu lăsați ca binele vostru să fie defăimat,

\VS{17}căci Împărăția lui Dumnezeu nu este mâncare și băutură, ci dreptate, pace și bucurie în Duhul Sfânt.

\VS{18}Căci cel ce slujește lui Hristos în aceste lucruri este plăcut lui Dumnezeu și aprobat de oameni.

\VS{19}Așadar, să urmăm lucrurile care duc la pace și lucrurile prin care ne putem edifica unii pe alții.

\VS{20}Nu răsturnați lucrarea lui Dumnezeu de dragul mâncării. Într-adevăr, toate lucrurile sunt curate, însă este rău pentru acel om care creează o piedică prin mâncare.

\VS{21}Este bine să nu mănânci carne, să nu bei vin și să nu faci nimic prin care fratele tău să se poticnească, să fie jignit sau să slăbească.


\par }{\PP \VS{22}Aveți credință? Ai credință în fața lui Dumnezeu. Fericit este cel care nu se judecă pe sine însuși în ceea ce aprobă.

\VS{23}Dar cel care se îndoiește este condamnat dacă mănâncă, pentru că nu este din credință; și tot ce nu este din credință este păcat.


\par }{\PP \VS{24}Și acum, Celui ce poate să vă întărească, după Buna Vestire a mea și după propovăduirea lui Isus Hristos, după descoperirea tainei, care a fost ținută ascunsă în veacuri îndelungate,

\VS{25}dar care acum s-a descoperit și, prin Scripturile profeților, după porunca Dumnezeului veșnic, a fost făcută cunoscută tuturor neamurilor, pentru ascultarea credinței,

\VS{26}pentru singurul și înțeleptul Dumnezeu, prin Isus Hristos, căruia I se cuvine slava în veci! Amin.



\par }\Chap{15}{\PP \VerseOne{1}Iar noi, care suntem puternici, trebuie să suportăm slăbiciunile celor slabi și să nu ne mulțumim pe noi înșine.

\VS{2}Fiecare dintre noi să mulțumească aproapelui său, pentru ceea ce este bine, ca să-l zidim.

\VS{3}Căci nici măcar Hristos nu s-a mulțumit pe sine însuși. Ci, după cum este scris: “Reproșurile celor care vă ocărau pe voi au căzut asupra mea”.

\VS{4}Căci tot ceea ce a fost scris mai înainte a fost scris pentru învățătura noastră, pentru ca, prin perseverență și prin încurajarea Scripturilor, să avem nădejde.

\VS{5}Or, Dumnezeul perseverenței și al încurajării să vă dea să fiți de aceeași părere unii cu alții, potrivit cu Hristos Isus,

\VS{6}pentru ca, cu o singură gură, să slăviți cu un singur gând pe Dumnezeul și Tatăl Domnului nostru Isus Hristos.


\par }{\PP \VS{7}De aceea, acceptați-vă unii pe alții, cum și Hristos v-a acceptat pe voi, spre slava lui Dumnezeu.

\VS{8}Acum spun că Hristos a fost făcut rob al circumciziei pentru adevărul lui Dumnezeu, ca să confirme promisiunile făcute părinților,

\VS{9}și ca neamurile să slăvească pe Dumnezeu pentru îndurarea Sa. După cum este scris,

\par }{\Q “De aceea te voi lăuda printre neamuri.

\par }{\QB și voi cânta în numele Tău.”


\par }{\PP \VS{10}El spune din nou,

\par }{\Q “Bucurați-vă, neamuri, împreună cu poporul Său”.


\par }{\PP \VS{11}Din nou,

\par }{\Q “Lăudați pe Domnul, voi toate neamurile!

\par }{\QB Toate popoarele să îl laude.”


\par }{\PP \VS{12}Isaia spune din nou,

\par }{\Q “Acolo va fi rădăcina lui Iesei,

\par }{\QB cel care se ridică să stăpânească peste neamuri;

\par }{\QB în El vor nădăjdui neamurile.”


\par }{\PP \VS{13}Și Dumnezeul nădejdii să vă umple de toată bucuria și pacea în credință, ca să sporiți în nădejde prin puterea Duhului Sfânt.


\par }{\PP \VS{14}Și eu însumi sunt încredințat de voi, frații mei, că voi înșivă sunteți plini de bunătate, plini de toată știința, în stare să îndemnați și pe alții.

\VS{15}Dar vă scriu cu mai multă îndrăzneală în parte ca să vă reamintesc, datorită harului care mi-a fost dat de Dumnezeu,

\VS{16}ca să fiu slujitor al lui Isus Hristos pentru neamuri, slujind ca preot al Bunei Vestiri a lui Dumnezeu, pentru ca jertfa neamurilor să fie acceptabilă, sfințită prin Duhul Sfânt.

\VS{17}Prin urmare, mă laud în Hristos Isus în ceea ce privește lucrurile care țin de Dumnezeu.

\VS{18}Căci nu îndrăznesc să vorbesc despre alte lucruri decât despre cele pe care Hristos le-a lucrat prin mine pentru ascultarea neamurilor, prin cuvânt și prin faptă,

\VS{19}în puterea semnelor și a minunilor, în puterea Duhului lui Dumnezeu; astfel încât, din Ierusalim și de jur împrejur până în Iliricum, am propovăduit pe deplin Vestea cea Bună a lui Hristos;

\VS{20}da, făcându-mi scopul de a propovădui Vestea cea Bună, nu acolo unde Hristos a fost deja numit, ca să nu construiesc pe temelia altuia.

\VS{21}Ci, după cum este scris,

\par }{\Q “Ei vor vedea, la care nu a venit nicio veste despre el.

\par }{\QB Cei care nu au auzit vor înțelege.”


\par }{\PP \VS{22}De aceea și eu am fost împiedicat de atâtea ori să vin la voi,

\VS{23}dar acum, neavând loc în aceste ținuturi și având de atâția ani dorința de a veni la voi,

\VS{24}voi veni la voi ori de câte ori voi merge în Spania. Căci sper să te văd în călătoria mea și să fiu ajutat de tine în drumul meu acolo, dacă mai întâi mă voi putea bucura de compania ta pentru o vreme.

\VS{25}Dar acum, zic, mă duc la Ierusalim, pentru a sluji sfinților.

\VS{26}Căci Macedonia și Ahaia au binevoit să facă o anumită contribuție pentru săracii dintre sfinții care sunt la Ierusalim.

\VS{27}Da, a fost pe placul lor, și ei sunt datori. Căci, dacă neamurile au fost făcute părtașe la lucrurile lor spirituale, le sunt datoare să le slujească și în cele materiale.

\VS{28}Așadar, după ce voi fi împlinit aceasta și le voi fi pecetluit acest rod, voi pleca pe drumul vostru spre Spania.

\VS{29}Știu că, atunci când voi veni la voi, voi veni în plinătatea binecuvântării Binecuvântării Bunei Vestiri a lui Hristos.


\par }{\PP \VS{30}Vă rog, fraților, pentru Domnul nostru Isus Hristos și pentru dragostea Duhului Sfânt, să vă rugați împreună cu mine în rugăciunile voastre către Dumnezeu pentru mine,

\VS{31}ca să fiu izbăvit de cei neascultători din Iudeea și ca slujba pe care o am pentru Ierusalim să fie plăcută sfinților,

\VS{32}ca să vin la voi cu bucurie, prin voia lui Dumnezeu, și împreună cu voi să găsesc odihnă.

\VS{33}Acum, Dumnezeul păcii să fie cu voi toți. Amin.



\par }\Chap{16}{\PP \VerseOne{1}Vă încredințez pe Febe, sora noastră, care este slujitoare a adunării din Cencreea,

\VS{2}ca să o primiți în Domnul, în chip vrednic de sfinți, și să o ajutați în tot ceea ce va avea nevoie de la voi, căci ea însăși a ajutat pe mulți și pe mine însumi.


\par }{\PP \VS{3}Salutați pe Prisca și pe Acuila, tovarășii mei de muncă în Hristos Isus,

\VS{4}care și-au riscat gâtul pentru viața mea și cărora nu numai eu le mulțumesc, ci și toate adunările neamurilor.

\VS{5}Salutați adunarea care se află în casa lor. Salutați-l pe Epaeneț, preaiubitul meu, care este cel dintâi rod al Ahaiei pentru Hristos.

\VS{6}Salutați-o pe Maria, care a muncit mult pentru noi.

\VS{7}Salutați-i pe Andronic și pe Junia, rudele mele și tovarășii mei de închisoare, care sunt remarcabili printre apostoli, care au fost și ei în Hristos înaintea mea.

\VS{8}Salutați-l pe Amplias, iubitul meu în Domnul.

\VS{9}Salutați pe Urbanus, tovarășul nostru de muncă în Hristos, și pe Stachys, iubitul meu.

\VS{10}Salutați pe Apelles, cel aprobat în Hristos. Salutați-i pe cei din casa lui Aristobulus.

\VS{11}Salutați pe Irodion, ruda mea. Salutați pe cei din casa lui Narcis, care sunt în Domnul.

\VS{12}Salutați pe Tryphaena și Tryphosa, care lucrează în Domnul. Salutați-o pe Persis, cea iubită, care a lucrat mult în Domnul.

\VS{13}Salutați pe Rufus, cel ales în Domnul, pe mama lui și pe ai mei.

\VS{14}Salutați pe Asincritus, Flegon, Hermes, Patrobas, Hermas și pe frații care sunt cu ei.

\VS{15}Salutați pe Filologus și pe Iulia, pe Nereus și pe sora lui, pe Olimpa și pe toți sfinții care sunt cu ei.

\VS{16}Salutați-vă unii pe alții cu o sărutare sfântă. Vă salută adunările lui Hristos.


\par }{\PP \VS{17}Acum, fraților, vă rog să vă uitați la cei ce fac dezbinări și prilejuri de poticnire, împotriva învățăturii pe care ați învățat-o, și să vă depărtați de ei.

\VS{18}Căci cei ca aceștia nu slujesc Domnului nostru Isus Hristos, ci propriei lor pântece; și prin vorbirea lor netedă și lingușitoare înșală inimile celor nevinovați.

\VS{19}Căci ascultarea voastră a devenit cunoscută de toți. De aceea mă bucur pentru voi. Dar doresc ca voi să fiți înțelepți în ceea ce este bun, dar nevinovați în ceea ce este rău.

\VS{20}Și Dumnezeul păcii îl va zdrobi repede pe Satana sub picioarele voastre.

\par }{\PP Harul Domnului nostru Iisus Hristos să fie cu voi.


\par }{\PP \VS{21}Timotei, tovarășul meu de muncă, vă salută, precum și Lucius, Iason și Sosipater, rudele mele.

\VS{22}Eu, Tertius, care scriu scrisoarea, vă salut în Domnul.

\VS{23}Gaius, gazda mea și gazda întregii adunări, vă salută. Erastus, trezorierul orașului, vă salută, la fel ca și fratele Quartus.

\VS{24}Harul Domnului nostru Isus Hristos să fie cu voi toți! Amin.

\VS{25}\FTNT{*}{{\FR 16:25 }{\FT{TR plasează Romani 14:24-26 la sfârșitul lui Romani în loc de la sfârșitul capitolului 14 și numără aceste versete 16:25-27.}}}


\par }